\documentclass[fontsize=14pt]{scrartcl}
\setkomafont{disposition}{\normalfont\bfseries}

% Standard American Math Society packages provide theorem environments, symbols,
% etc.
\usepackage{amsthm}
\usepackage{amsmath}
\usepackage{amssymb}

%margins
\usepackage[margin=.75in]{geometry}

%links table of contents
\usepackage{hyperref}
\hypersetup{
    colorlinks,
    citecolor=black,
    filecolor=black,
    linkcolor=blue,
    urlcolor=black
}

% 
\usepackage{cleveref}

%header
\usepackage{fancyhdr}
\setlength{\footskip}{15pt}
\pagestyle{fancy}

%enumerate
\usepackage{enumitem}


\title{Notes}
\author{Nicholas Vantine}
\date{}

\input{macros}
% Define theorem, definition, lemma, etc. to share the same counter
\newtheorem{theorem}{Theorem}[section]
\newtheorem{lemma}[theorem]{Lemma}
\newtheorem{corollary}[theorem]{Corollary}
\newtheoremstyle{defnstyle}
  {\topsep}   % Space above
  {\topsep}   % Space below
  {\normalfont}  % Body font
  {0pt}         % Indent amount
  {\bfseries}   % Theorem head font
  {.}          % Punctuation after theorem head
  {.5em}     % Space after theorem head
  {}          % Theorem head spec (empty for plain)
\theoremstyle{defnstyle}
\newtheorem{definition}[theorem]{Definition}
\theoremstyle{definition}
\newtheorem{example}[theorem]{Example}
\newtheorem{exercise}[theorem]{Exercise}


\setmonofont{DejaVu Sans Mono} % has broad unicode coverage including ℝ
\tcbuselibrary{skins, breakable}

\newcounter{leancode}
\counterwithin{leancode}{section}

\newtcolorbox{leanbox}[2][]{
    colback=gray!5,
    colframe=gray!50,
    coltitle=black,
    fonttitle=\bfseries\small,
    breakable,
    boxrule=0.5pt,
    arc=2pt,
    title={Lean Code \refstepcounter{leancode}\theleancode: #2},
    #1
}
\begin{document}

\maketitle
\newpage
\tableofcontents
\bibliographystyle{plain}


\newpage


\section{Foundations}


A relation simply relates elements in two different sets by delineating a subset of related pairs in the Cartesian product.

\begin{definition}
Let $X$ and $Y$ be sets. A \emph{relation} over $X \times Y$ is a subset $R \subseteq X \times Y$. If $Y = X$ then it is said to be \emph{homogeneous}, and $R$ is simply said to be a relation over $X$. If $(x,y) \in R$ then we say $x$ is related to $y$ and write $xRy$.
\end{definition}

\begin{example}
Let $X$ and $Y$ be sets. Then the following are relations:
\begin{enumerate}
    \item The \emph{discrete relation}, $X \times Y$.
    \item The \emph{empty relation}, $\emptyset \subseteq X \times Y$.
    \item The \emph{diagonal relation}, $\{(x,x) : x \in X\}$.
\end{enumerate}
\end{example}

\subsection*{Exercises}

\begin{exercise}
Give a necessary and sufficient condition for the empty, discrete, and diagonal relations to be equal.
\end{exercise}

\begin{proof}[Solution]
    Let $Y = X$ and $X = \emptyset$
\end{proof}

\begin{exercise}
If $(X_1, R_1)$ and $(X_2, R_2)$ are pairs of sets and relations on them, then the rule,
\[
(x,y) \, R_1 \times R_2 \, (x', y') \quad \text{iff} \quad xR_1x' \text{ and } yR_2y',
\]
defines a relation $R_1 \times R_2$ on $X_1 \times X_2$. Express this relation as a subset of $(X_1 \times X_2)^2$.
\end{exercise}

\begin{proof}[Solution]
    \[
        R_1 \times R_2 = \{ ((x,y),(x',y')) \in X_1 \times X_2 : xR_1x' \text{ and } yR_2y' \}
    \]
\end{proof}

\newpage


\begin{definition}
     Let $X$ be a set and $R$ a relation over $X.$ We say $R$ is
     \begin{enumerate}
         \item \emph{reflexive} if $xRx$,  $\forall x \in X$,
         \item \emph{symmetric} if $xRy$ implies $yRx$,  $\forall x,y \in X$,
         \item \emph{transitive} if $xRy$ and $yRz$ implies $xRz$, $\forall x,y,z \in X$,
         \item is an \emph{equivalence relation} on $X$ if it is reflexive, symmetric, and transitive.
     \end{enumerate}
\end{definition}

\textbf{Notation.} An arbitrary equivalence relation is typically denoted $\sim$.


\newpage


\section{Advanced Calculus \cite{Abbott2015}}


\begin{exercise}
    Show 
    \begin{align*}
        \left(\bigcup_{i=1}^{\infty} A_i\right)^c = \bigcap_{i=1}^{\infty} A_i^c.
    \end{align*}
\end{exercise}

\begin{proof}
    We show both inclusions via a chain of biconditionals.

    Let $x$ be arbitrary. Then
    \begin{align*}
        x \in \left(\bigcup_{i=1}^{\infty} A_i\right)^c
        &\iff x \notin \bigcup_{i=1}^{\infty} A_i \\
        &\iff x \notin A_i \text{ for all } i \in \mathbb{N} \\
        &\iff x \in A_i^c \text{ for all } i \in \mathbb{N} \\
        &\iff x \in \bigcap_{i=1}^{\infty} A_i^c.
    \end{align*}

    Since each step above is a biconditional, the chain holds in both directions. Hence
    \begin{align*}
        \left(\bigcup_{i=1}^{\infty} A_i\right)^c = \bigcap_{i=1}^{\infty} A_i^c.
    \end{align*}
\end{proof}


\newpage


\section{Linear Algebra \cite{Axler2024}}

\begin{definition}[Addition]\label{def:addition}
    A binary operation $+: S \times S \to S$, $(s,t) \mapsto s+t$, for $s,t \in S$.
\end{definition}

\begin{definition}[Scalar Multiplication]\label{def:scalar-mult}
    A map $*: \F \times S \to S$, $(\lambda, s) \mapsto \lambda s$, for $\lambda \in \F,\ s \in S$.
\end{definition}



\begin{definition}
    A Vector Space $V$ is a set equipped with an \emph{addition} and \emph{scalar multiplication} that hold the following properties:
    \begin{itemize}
        \item \emph{Commutativity}: $\quad \forall u,v \in V, \quad u+v=v+u.$ 
        \item \emph{Associativity}: $\quad \forall u,v,w \in V, \quad (u + v) + w = u + (v + w).$ 
        \item \emph{Additive Identity}: $\quad \exists 0 \in V, \forall v \in V, \quad v + 0 = v.$
        \item \emph{Additive Inverse}: $\quad \forall v \in V, \exists w \in V, \quad v + w = 0.$
        \item \emph{Multiplicative Identity}: $\quad \forall v \in V, \quad 1v = v.$
        \item \emph{Distributive Laws}: $\quad \forall u,v \in V$ and $ \forall s,t \in \F,$ the following hold:
            \[
            \begin{aligned}
                a(v + w) = av + aw \\ 
                (a + b)v = av + bv.
            \end{aligned}
            \]
    \end{itemize}
\end{definition}


\begin{theorem}
    $\forall v \in V, \exists! w \in V, \quad v + w = 0.$
\end{theorem}
\begin{proof}
    Suppose $V$ is a vector space. Let $ v \in V.$ Suppose $w$ and $w'$ are additive inverses of $v.$ Then
    \[
        w = w + 0 = w + (v +w') = (w + v) + w' = 0 + w' = w'.
    \]
    Thus $w = w'$. 
\end{proof}


\textbf{Remark.} We just proved a statement of the form,
\[
    \forall v \in V, [\exists u \in V, v + u = 0] \land [\forall w,w' \in V, ((v + w = 0) \land (v + w' = 0))\implies w' = w].
\]

While this makes the proof more complicated, I think it is important to see what is going on under the hood to truly understand the mechanisms that make this proof work. First we take the for all part of the statement, which says 'given any vector can you give me at least one of its inverses?' Surely, by the definition of a vector space a vector $u$ exists. Therefore we have the first part of conjunction completed. Now for the second part of the conjunction we have to assume that we are given 2 arbitrary inverses of $v$, and with that we can other other assumptions about vector spaces to always conclude that these 2 arbitrary elements are actually the same thing. Looking back at the proof you can see that the first part was not explicitly stated and proved at all but notice by working on the second part of the conjunction you are implicitly proving the existence part of the conjunctive statement.



\newpage

\begin{exercise}
    Prove that $-(-v) = v$ for every $v \in V.$
\end{exercise}

\begin{proof}
    Suppose $v \in V$, where $V$ is a Vector Space. 
    By definition $-v$ is the additive inverse of $v$, 
    \[
        v + (-v) = 0
    \]
    Also by definition, the additive inverse of $-v$ is, 
    \[
        -v + (-(v)) = 0
    \]
    Using commutativity we can rearrange one expression to see, 
    \[
        \begin{aligned}
            (-v) + v = 0 \\
            (-v) + (-(-v)) = 0
        \end{aligned}
    \]
    Now using the fact additive inverses are unique, we can conclude that 
    \[
        -(-v) = v.
    \]
\end{proof}

-- \href{https://live.lean-lang.org/#codez=JYWwDg9gTgLgBAWQIYwBYBtgCMBQOBuSUwSW6ApnAN4BqcAXHACoCeY5AVAL5wDaAggBNBAYQggQAcSgQArmDg0AunwQRBsinEC4hIqV5yADyTgtACnwNFASisBaM3fy2AvHEv03WFnjhxUSPiUqACMVpYA1HCOznBuAAwMbkjCAPoAduQA5qkAxkjpueTo7r7%2BgcEATFYxtlGOta5wiZ5wKYIZ2XkFRSW1ZVAA7nzteeIgKij%2BIWV2dtP2ke5xzbPzqNWMTnBRDk5N8WUBQUltaRQAZjDdhcXRoQB0MFAFAM7%2BlQ%2BvLBLWZUZIXLwNDAV5fH4gIA}{Click here for the computer verified proof.}



\newpage

\section{Real Analysis}


\section{Topology}


\section{Differential Geometry}

\newpage


\section{Convex Optimization \cite{cvxbook}}


\subsection{Convex Sets}


\begin{definition}
    A set $C$ is convex if the line segment between any two points in $C$ lies in $C$, i.e., 
    if for any $x_1,$ $x_2$ $\in$ $C$ and any $\theta$ with $ 0 \le \theta \le 1,$ we have
    \[
        \theta x_1 + (1 - \theta) x_2 \in C.
    \]
\end{definition}


\subsection*{Excercises}


\begin{exercise}
    Let $C \subseteq \R^n$ be a convex set, with $x_1, \ldots, x_k \in C$, and let $\theta_1, \ldots, \theta_k \in \R$ satisfy $\theta_i \geq 0$, $\theta_1 + \cdots + \theta_k = 1$. Show that $\theta_1 x_1 + \cdots + \theta_k x_k \in C$. (The definition of convexity is that this holds for $k = 2$; you must show it for arbitrary $k$.) \textit{Hint.} Use induction on $k$.
\end{exercise}

\begin{proof}
    We prove the result by induction on the number of points.

    For the base case, let $k = 2$. If $x_1, x_2 \in C$ and
    $\theta \in [0,1]$, then by the definition of convexity,
    \begin{equation*}
        (1-\theta)x_1 + \theta x_2 \in C.
    \end{equation*}
    Hence every convex combination of two points in $C$ belongs to $C$.

    Now suppose every convex combination of $\ell$ points in $C$
    belongs to $C$. That is, whenever
    \begin{equation*}
        x_i \in C,\qquad
        \theta_i \geq 0,\qquad
        \sum_{i=1}^{\ell}\theta_i = 1,
    \end{equation*}
    we have
    \begin{equation*}
        \sum_{i=1}^{\ell}\theta_i x_i \in C.
    \end{equation*}

    Now let $x_1,\ldots,x_{\ell+1}\in C$ and let
    $\phi_i \geq 0$ satisfy
    \begin{equation*}
        \sum_{i=1}^{\ell+1}\phi_i = 1.
    \end{equation*}
    We wish to show that
    \begin{equation*}
        \sum_{i=1}^{\ell+1}\phi_i x_i \in C.
    \end{equation*}

    If $\phi_{\ell+1}=1$, then
    $\phi_1=\cdots=\phi_\ell=0$, and therefore
    \begin{equation*}
        \sum_{i=1}^{\ell+1}\phi_i x_i
        =
        x_{\ell+1}
        \in C.
    \end{equation*}

    Now suppose $\phi_{\ell+1}<1$, and define
    \begin{equation*}
        z
        =
        \frac{1}{1-\phi_{\ell+1}}
        \sum_{i=1}^{\ell}\phi_i x_i.
    \end{equation*}

    Since
    \begin{equation*}
        \sum_{i=1}^{\ell}\phi_i
        =
        1-\phi_{\ell+1},
    \end{equation*}
    it follows that
    \begin{equation*}
        \sum_{i=1}^{\ell}
        \frac{\phi_i}{1-\phi_{\ell+1}}
        =
        1,
    \end{equation*}
    and
    \begin{equation*}
        \frac{\phi_i}{1-\phi_{\ell+1}}
        \geq 0,
        \qquad i=1,\ldots,\ell.
    \end{equation*}

    Therefore, the induction hypothesis implies that
    \begin{equation*}
        z \in C.
    \end{equation*}

    Finally,
    \begin{equation*}
    \begin{aligned}
        \sum_{i=1}^{\ell+1}\phi_i x_i
        &=
        \sum_{i=1}^{\ell}\phi_i x_i
        + \phi_{\ell+1}x_{\ell+1} \\
        &=
        (1-\phi_{\ell+1})z
        + \phi_{\ell+1}x_{\ell+1}.
    \end{aligned}
    \end{equation*} \\

    Since $z,x_{\ell+1}\in C$ and
    $\phi_{\ell+1}\in[0,1]$, convexity of $C$ yields
    \begin{equation*}
        (1-\phi_{\ell+1})z
        + \phi_{\ell+1}x_{\ell+1}
        \in C.
    \end{equation*}

    Hence
    \begin{equation*}
        \sum_{i=1}^{\ell+1}\phi_i x_i
        \in C.
    \end{equation*}

    Therefore, by induction, every finite convex combination of points in $C$
    belongs to $C$.
\end{proof}


\newpage


\begin{exercise}
    What is the distance between the parallel hyperplanes
    \begin{align*}
        \{x \in \R^n \mid a^T x = b_1\}
        \quad \text{and} \quad
        \{x \in \R^n \mid a^T x = b_2\}\, ?
    \end{align*}
\end{exercise}

\begin{proof}
    Let
    \[
        H_1=\{x\in\R^n\mid a^Tx=b_1\}
        \qquad\text{and}\qquad
        H_2=\{x\in\R^n\mid a^Tx=b_2\}.
    \]

    Choose arbitrary points $x_1\in H_1$ and $x_2\in H_2$. By the
    orthogonal decomposition theorem, there exist $t\in\R$ and
    $y\in\R^n$ such that
    \[
        x_1-x_2=ta+y,
        \qquad y\perp a.
    \]

    In particular,
    \[
        y=x_1-x_2-ta.
    \]

    Since $ta$ is parallel to $a$, it is perpendicular to both
    hyperplanes. Therefore, the distance between the hyperplanes is
    the norm of this component:
    \[
        d=\|ta\|.
    \]

    Since $y\perp a$, we also have $y\perp ta$. Hence,
    \[
    \begin{aligned}
        0
        &=\langle ta,x_1-x_2-ta\rangle \\
        &=\langle ta,x_1\rangle
          -\langle ta,x_2\rangle
          -\langle ta,ta\rangle \\
        &=t\langle a,x_1\rangle
          -t\langle a,x_2\rangle
          -t^2\|a\|^2 \\
        &=tb_1-tb_2-t^2\|a\|^2.
    \end{aligned}
    \]

    If $b_1=b_2$, then the two hyperplanes are equal and $d=0$.
    Now suppose $b_1\neq b_2$. Then $t\neq0$, so
    \[
        t\|a\|^2=b_1-b_2.
    \]

    Thus,
    \[
        t=\frac{b_1-b_2}{\|a\|^2}.
    \]

    Therefore,
    \[
    \begin{aligned}
        d
        &=\|ta\| \\
        &=\left\|
            \frac{b_1-b_2}{\|a\|^2}a
        \right\| \\
        &=\frac{|b_1-b_2|}{\|a\|^2}\|a\| \\
        &=\frac{|b_1-b_2|}{\|a\|}.
    \end{aligned}
    \]

    Hence, the distance between the two parallel hyperplanes is
    \[
        \frac{|b_1-b_2|}{\|a\|}.
    \]
\end{proof}


\newpage


\begin{exercise}
    \textit{Voronoi description of halfspace.} Let $a$ and $b$ be distinct points in $\mathbf{R}^n$. Show that the set of all points that are closer (in Euclidean norm) to $a$ than $b$, i.e.,
        $
        \{x \mid \|x - a\|_2 \leq \|x - b\|_2\},
        $
    is a halfspace. Describe it explicitly as an inequality of the form $c^T x \leq d$. Draw a picture.
\end{exercise}

\begin{proof}
    Since both sides are nonnegative,
    \[
        \|x-a\|_2 \le \|x-b\|_2 \iff \|x-a\|_2^2 \le \|x-b\|_2^2.
    \]
    Expanding,
    \[
        \|x-a\|_2^2 = x^Tx - 2a^Tx + a^Ta, \qquad \|x-b\|_2^2 = x^Tx - 2b^Tx + b^Tb.
    \]
    The inequality becomes
    \[
        x^Tx - 2a^Tx + a^Ta \le x^Tx - 2b^Tx + b^Tb,
    \]
        and the quadratic terms $x^Tx$ cancel, leaving
    \[
        2(b-a)^Tx \le b^Tb - a^Ta.
    \]
    Thus
    \[
        \{x \mid \|x-a\|_2 \le \|x-b\|_2\} = \{x \mid c^Tx \le d\}, \qquad
        c = b-a, \quad d = \tfrac{1}{2}\big(\|b\|_2^2 - \|a\|_2^2\big).
    \]
    Since $a \ne b$, $c \ne 0$, so this is indeed a halfspace. Geometrically, it is the halfspace bounded by the perpendicular bisector of the segment $[a,b]$, containing $a$, with boundary normal vector $c = b - a$ pointing from $a$ toward $b$.
\end{proof}


\newpage


\bibliography{ref}

\end{document}